\chapter{Simulation Validation}\label{ch:validation}

This chapter reports the tests that decide whether the INDI controllers are accepted.
Each criterion of \cref{tab:criteria} is a \emph{gate}: a test with fixed thresholds.
Gate G$n$ is criterion C$n$.
Six gates are \passed{}: G1, G2, G4, G5, G7 and G8.
Two are \notpassed{}: G3 (robustness) and G6 (defaults and identification, which needs G3).
All evidence comes from simulation, on one reference airframe: a 5-inch quadcopter of about 600~g.
No real flight was made.
\Cref{ch:limits} states what this means.

\section{Method}\label{sec:val_method}

\subsection{Gates as test definitions}\label{sec:val_gates}

Each gate was written as a test definition before its first run: what is flown, what is measured and which threshold applies.
The thresholds do not depend on the controller that is being judged.
In particular, a PID controller was flown on the same scenarios as a reference.
Its results are shown beside the INDI results where they exist.
They are never a pass condition.
A gate passes only when its own thresholds are met.

A verdict holds for one firmware.
The gates G2, G4 and G5 were flown on the final configuration.
If the firmware changes later, the rows that the change affects must be flown again.
The robustness re-flights used a simulator build in which the motor-fault injector was repaired.
Earlier robustness runs with several simultaneous motor limits were declared invalid, because the injector had applied the limits about one second apart.

\subsection{Two scoring scripts per metric}\label{sec:val_keys}

Every important metric is computed twice, by two scoring scripts written independently.
The report calls them key~1 and key~2.
A result counts when both keys agree.
When they disagree, the case is reported and the gate definition decides, not the script that happens to be more favourable.
The scoring scripts themselves are tested.
Each one has a self-test that must fail when the check is deliberately broken.

\subsection{Ground truth, seeds and run provenance}\label{sec:val_truth}

The simulator publishes the true state of the vehicle, stamped in simulation time.
All pass conditions are judged on this truth, not on the estimator output.
Each test is flown with several random seeds, three for most gates.
Every run records the code commit, whether the source tree was clean, the hash of the simulator binary, and the hashes of the simulation models and of the parameters.
The run script refuses to start with a modified source tree or an out-of-date binary.
Most INDI flights ran at four times real time.
The flights that decide the robustness gate were repeated at real time (\cref{sec:val_g3}), and the limits chapter discusses the remaining risk.

\subsection{Owner amendments}\label{sec:val_amend}

After the first run of a test, the project owner amended the test where it had a defect.
Each change was logged before the results were re-scored.
Where a gate result below depends on an amendment, the table says so.
\Cref{sec:val_amendments} lists the amendments in one place.

\section{G1: evidence integrity}\label{sec:val_g1}

G1 does not test the controller.
It tests whether the evidence can be trusted.
\Cref{tab:val_g1} lists its seven items.

\begin{table}[tbp]
\caption{G1 items and results.}\label{tab:val_g1}
\centering
\small
\begin{tabularx}{\textwidth}{@{}l L{4.9cm} Y@{}}
\toprule
Item & Test & Result \\
\midrule
G1.1 & Logger completeness: every published sample is logged, at most 0.1~\% missing & Met. Logger shortfall within the limit. Simulator holes: 0.012~\% of steps, single holes at most 6.3~ms. These fail the original gap rule, which the owner replaced by this one \\
G1.2 & True state is stamped in simulation time & Met \\
G1.3 & Two scoring keys, with the attitude angle computed the same way in both & Met \\
G1.4 & Scoring self-tests fail when the check is broken & Met \\
G1.5 & Each fault injector and plant variant is checked before its first campaign & Deferred to each injector by design; not a failed item \\
G1.6 & Mode-change timing: spread at most 10~ms, latency at most 20~ms & Met: spread within one commander period (10~ms), latency within two \\
G1.7 & Run provenance recorded, dirty tree or stale binary refused & Met. The refusal was tested by breaking each condition: dirty tree, stale binary, newer header, missing input, missing binary \\
\bottomrule
\end{tabularx}
\end{table}

\begin{keyresult}[G1 verdict: \passed]
The logs are complete enough and the scoring is checked.
G1 carries five evidence limits into every later verdict.
\begin{itemize}
  \item Logs have two kinds of holes (simulator and controller). A run with more than 0.1~\% holes is flown again; it is not a failure.
  \item Analysis windows anchor on the logged mode change, which is timed to within 10~ms.
  \item The estimated height differs from the true height by up to 0.15~m in hover, and the position-hold spread is $\pm 0.15$~m.
  \item The estimated heading is off by up to 1.6~deg, constant in a hover. No yaw threshold below 1.6~deg is judged on truth.
  \item The bench parameter sets switch off the barometer auto-calibration.
\end{itemize}
\end{keyresult}

\section{G2: stability margins}\label{sec:val_g2}

G2 measures the stability margins of every loop in flight, with a multisine injected at a loop-break point (\cref{ch:analysis}).
It then checks attitude steps.
\Cref{tab:val_g2_tests} gives the thresholds.

\begin{table}[tbp]
\caption{G2 thresholds.}\label{tab:val_g2_tests}
\centering
\small
\begin{tabularx}{\textwidth}{@{}L{4.2cm} Y@{}}
\toprule
Quantity & Threshold \\
\midrule
Phase margin & At least 45~deg \\
Gain margin & At least 6~dB; no limit where the phase never reaches $-180^\circ$ in 0.2 to 8~Hz \\
Gain reduction near $-180^\circ$ & Where $|L|>1$ within 10~deg of $-180^\circ$, the value $20\log_{10}|L|$ must also be at least 6~dB \\
Delay margin & At least two control periods plus telemetry latency (never the binding limit) \\
Key agreement & Phase margins within 3~deg, gain margins within 1~dB, and both keys find a gain margin or neither does \\
Holes & The hole check of G1 passes in every run \\
Step response, every seed & Roll and pitch: overshoot at most 5~\%, 2~\% settling at most 450~ms. Yaw: overshoot at most 5~\%, settling at most 1350~ms \\
\bottomrule
\end{tabularx}
\end{table}

The step thresholds are absolute.
They were derived from a PID reference flown beforehand on three seeds: roll and pitch overshoot 0~\% and settling at most 412~ms, yaw overshoot at most 0.3~\% and settling at most 1240~ms.

\paragraph{Flights.}
The outer loops were flown in 24 cells: three axes, two break points, two controller rates (400 and 800~Hz) and two seeds.
Two spot cells cover the inner loops: the roll rate loop at 800~Hz and the pitch attitude loop at 400~Hz.
The step test used three seeds on roll, pitch and yaw.
The 250~Hz controller rate was not measured.

\begin{table}[tbp]
\caption{G2 results at the final configuration. Worst case over all flown cells.}\label{tab:val_g2}
\centering
\small
\begin{tabularx}{\textwidth}{@{}L{4.6cm} L{3.2cm} L{2.2cm} Y@{}}
\toprule
Quantity & Worst case & Threshold & Result \\
\midrule
Outer-loop phase margin & 49.1~deg & 45~deg & 24 of 24 meet it \\
Outer-loop gain margin, horizontal axes & 13.5~dB & 6~dB & 24 of 24 meet it \\
Plateau gain margin (second key) & 12.88~dB & 6~dB & Met \\
Outer-loop delay margin & 229~ms & not binding & Met \\
Hole check & 24 of 24 & pass & One cell failed and was flown again \\
Agreement of the two keys & 21 of 24 cells & agree & Three cells blocked (see text) \\
Rate loop, roll spot & PM 46.9~deg, GM 15.7~dB, DM 75.5~ms & PM 45~deg & Met, thin \\
Attitude loop, pitch spot & PM 48.0~deg, GM 13.7~dB at 4.29~Hz, DM 93.5~ms & PM 45~deg & Met, thin \\
Steps, roll and pitch & overshoot 0.65~\%, settling 347.5~ms & 5~\%, 450~ms & 9 of 9 pass, both keys \\
Steps, yaw & overshoot at most 0.17~\%, settling 1196.25~ms & 5~\%, 1350~ms & 89~\% of the limit \\
\bottomrule
\end{tabularx}
\end{table}

\paragraph{Reading the result.}
The worst hole-check cell had one setpoint gap of 22.5~ms against a 10~ms limit.
It was flown again and passed with a phase margin of 56.9~deg, a gain margin of 13.9~dB and a delay margin of 0.485~s.
The earlier run of this cell had a phase margin of 59.59~deg and a gain margin of 13.83~dB.
Three outer cells are blocked on key agreement only; both keys are above the limits in each of them:
\begin{itemize}
  \item one velocity cell: phase margin 71.3~deg on key~1 and 67.8~deg on key~2 (difference $-3.48$~deg, limit 3~deg);
  \item one acceleration cell on the vertical axis: key~2 finds a gain margin of 16.3~dB, key~1 finds none;
  \item one acceleration cell on the vertical axis: gain margin 22.4~dB on key~1 and 17.8~dB on key~2 (difference $-4.6$~dB).
\end{itemize}
These cells were not flown again.
In the primary scorer, the vertical outer loop has no $-180^\circ$ crossing, hence no gain margin, in 7 of its 8 cells.

The attitude spot cell has a second $-180^\circ$ crossing at 0.237~Hz with $|L|=7.3$.
It is read as a gain-reduction margin of 17.3~dB.
A literal signed reading would print $-17.31$~dB and fail.
The gate verdict uses the first reading, which the owner's amendment prescribes.
The rate and attitude spot margins exceed the limit by only 1.9 and 3.0~deg.
Pitch and yaw rate margins were not flown at the final configuration.

\indifigure{step_attitude}{Attitude step responses of the INDI controller and of the PID reference in simulated hover (15~deg roll and pitch, 45~deg yaw). Both settle inside the 2~\% band within the acceptance limits (450~ms, yaw 1350~ms). INDI overshoots by less than 1~\%; the PID reference barely overshoots.}{fig:val_steps}

\begin{keyresult}[G2 verdict: \passed]
All 24 outer-loop cells and both spot cells meet the margins, and all 9 steps pass on both keys.
Evidence limits: three cells are blocked on key agreement only; the spot margins are thin; there is no second key on the spot cells; 250~Hz was not measured.
\end{keyresult}

\section{G3: robustness}\label{sec:val_g3}

\subsection{The test}

G3 is a campaign of 300 valid simulated flights of 100~s.
240 flights are Latin-hypercube draws over the factors of \cref{tab:val_g3_factors}, and 60 are corner cases (six named corners C1 to C6, ten flights each) that combine three or four extreme factors.
Each flight is: arm, take-off, hover, roll steps, a lemniscate trajectory, then an event hover with a fault or a gust, and a hold.
Seventy percent of the flights use position mode, 20~\% attitude mode and 10~\% rate mode with zero sticks.
The noise of G4 (typical profile) is always on.

\begin{table}[tbp]
\caption{Factors drawn in the robustness campaign.}\label{tab:val_g3_factors}
\centering
\small
\begin{tabularx}{\textwidth}{@{}L{5.2cm} Y@{}}
\toprule
Factor & Range \\
\midrule
Controller gain versus plant gain, per axis & Owner box 0.8 to 1.2 times (see text) \\
Motor time constant & 0.5 to 3 times nominal, with different rise and fall values \\
Added loop delay & 0 to 10~ms \\
Mass; inertia & $\pm 30$~\%; $\pm 30$~\% per axis \\
Available thrust (battery sag) & 1.0 down to 0.7, in 40~\% of flights \\
Weak motor & one motor at 1.0 to 0.7, in 30~\% of flights \\
Saturation corner & thrust down to twice the hover thrust, 10~\% of flights \\
Telemetry & 0 to 3 frames latency; dropouts, loss and frozen values \\
Vibration & 1 to 5 times a base level \\
GNSS latency & 80 to 250~ms, with a mismatch of up to $\pm 50$~ms in the estimator setting \\
Wind & steady 0 to 5~m/s plus a gust with peak 3 to 9~m/s \\
\bottomrule
\end{tabularx}
\end{table}

Six items are checked on every flight, on the true state:
\begin{enumerate}
  \item no loss of control;
  \item no sustained oscillation of the rates or the position;
  \item steady state in the last 10~s of a hold of at least 20~s: horizontal RMS error at most 0.15~m in calm air and 0.30~m with wind or gust, height offset at most 0.10~m (0.20~m with sag or a weak motor);
  \item tracking: lemniscate horizontal RMS at most 0.5~m, vertical RMS at most 0.3~m;
  \item recovery: after each event the vehicle is back within 0.3~m and 3~deg in at most 5~s and does not leave the 0.3~m band for the next 5~s;
  \item rate mode: body-rate RMS at most 5~deg/s over the last 20~s with zero sticks.
\end{enumerate}
The campaign passes only if there is no loss of control, no sustained oscillation, and items 3 to 6 are met in every valid run.
There is no quota of allowed failures.
PID was flown on the same draws as a reference.

\subsection{Results}

No attempt passed.
\Cref{tab:val_g3} lists the four attempts.
The first is the full campaign, flown on the final firmware.
The others tried to repair the failures and stopped at the first confirmed failure.

\begin{table}[tbp]
\caption{Attempts at the robustness gate. All are NOT PASSED.}\label{tab:val_g3}
\centering
\small
\begin{tabularx}{\textwidth}{@{}L{3.2cm} Y@{}}
\toprule
Attempt & Result \\
\midrule
Full campaign, 300 of 300 flown, final firmware &
  Item 1 fails: one 2~Hz growing pitch oscillation ends in a fall (corner C1), and one flight tilts past 60~deg at the fault (corner C3, saturation plus a weak motor). Item 2 fails on 79 of 297 scored runs. Item 3 passes in 94.95~\% (15 misses). \\
Re-fly with the controller time constant drawn like an identification error, halted after 25 runs &
  One oscillation failure was not reproduced at real time and was not counted. The saturation draw lost control again, at 60.78~s at four times real time and 61.03~s at real time. \\
Three repair candidates in the rate loop, each on its own build &
  All missed the pass line on both keys (details below). \\
Last route, 300 flights planned, halted at the third &
  Slow-motor draw: item 2 fails on both keys, and tracking (item 4) fails on draws 2 and 3. \\
\bottomrule
\end{tabularx}
\end{table}

\paragraph{Item 2, 79 of 297.}
The 79 runs are the runs in which at least one key found a sustained oscillation: key~1 found it in 73 runs, key~2 in 45, and 39 runs were found by both.
They are four different things.
Eight runs show a 10.75 to 11.5~Hz pitch line (high delay, fast motor, controller gain below the plant gain).
Thirty-seven runs show a 2.0 to 2.5~Hz line (high delay and slow motor); the corners C1 and C6 fail in all 19 scored runs.
Twenty-five runs are landing windows that only key~1 flags; key~2 skips landings by design.
Nine runs are borderline 2~Hz lines flagged by key~1 only.
Two of the 2~Hz and 11~Hz cases were repeated at real time and reproduced, and PID passes them.
One of the 11~Hz draws also fails with PID, on another axis.
The other 77 flagged runs were not repeated at real time; the verdict does not depend on them.

\paragraph{Two failures that are specific to INDI.}
\begin{itemize}
  \item \emph{Slow motors.} On one draw the motors are 2.4 to 2.8 times slower than nominal. The INDI pitch rate oscillates at 2.00~Hz, sustained, on both keys, and again at real time. The local power peak (2.0 to 2.25~Hz) is about 30~dB above PID. PID on the same draw and seed passes.
  \item \emph{Saturation with a weak motor.} In corner C3 the thrust scale is 0.351, motor 3 is at 0.709 and the mass is 30~\% above nominal. INDI tilts past 60~deg at the fault: at 60.67~s (four times real time) and 60.91~s (real time) with the actuator time constant at 0.033~s, and again at 61.03~s with a drawn value of 0.0177~s. PID flies the same draw without loss of control; its tilt stays below 40~deg. A different time constant does not change the outcome.
\end{itemize}

\paragraph{Tracking, item 4.}
On three draws the lemniscate error exceeds 0.5~m (INDI against PID: 0.418 and 0.768~m, 0.834 and 0.862~m, 0.553 and 0.917~m).
The owner ruled that a row on which PID also fails the same item on the same draw and seed is a test defect, provided INDI is not worse.
These rows are reported and not counted against INDI.
Each has one flight per stack, so there is no spread estimate.

\paragraph{Repair candidates.}
All were rate-loop changes that were not adopted.
The first lowered the blend time constant between measured and modelled torque to 0.10~s: it fails items 1 and 2 on the saturation draw and item 2 on a second draw.
The second added a per-motor correction with a time constant of 0.2~s: it loses control on a joint draw and on the saturation draw.
The third added a bounded rate integrator: it loses control on the saturation draw and fails item 2 on the second draw.
Only 25 flights of the second attempt and 3 of the last were flown, so no firmware after the final configuration has a 300-flight result.

\paragraph{Recovery after events (item 5).}
The owner fixed item 5 to require two conditions: recovery within 5~s without a new excursion larger than 0.3~m, and, from 5 to 10~s after the event, position error at most 0.3~m and tilt error at most 3~deg.
A re-score of four pre-test runs reproduced the owner's expected outcome in all cases (six of six events).

\paragraph{Gain range as drawn.}
The recorded draws of the first campaign show controller gain factors between 0.62 and 1.69 times, wider than the amended box of \cref{tab:val_g3_factors}.
This was not resolved.
A narrower box would not change a verdict that rests on item 1 and item 2.

\indifigure{g3_campaign}{Robustness campaign: 300 simulated flights with randomly drawn plant and disturbance factors. 218 runs pass, 79 fail at least one item and 2 end in loss of control (one further run was not scored). Failures are more frequent with slow motors and large added delay. The campaign is not passed.}{fig:val_g3}

\begin{keyresult}[G3 verdict: \notpassed, final]
The owner closed the gate as not passed and stopped further attempts.
Two failures are INDI-specific: a 2~Hz pitch oscillation on slow motors, and loss of control with saturation plus a weak motor.
A cause was not found; \cref{sec:lim_g3} gives the one hypothesis and its status.
\end{keyresult}

\section{G4: noise and effort}\label{sec:val_g4}

G4 checks that sensor noise does not drive the motors too hard.
Three noise profiles (quiet, typical, harsh) were flown in hover and in cruise, three seeds each: 18 INDI flights.
Four metrics are measured on the motors.
R20 is the rotor-speed ripple RMS above 20~Hz, R4 the ripple RMS around 4~Hz, C20 the command RMS above 20~Hz, and the clip fraction is the share of samples at command 0 or 1.
\Cref{tab:val_g4_lim} gives the hover limits.
In cruise, R20 and C20 are allowed 1.5 times the hover limit.
Two further rows hold on the harsh profile in hover: G3 items 1 to 3, and no sustained 20 to 100~Hz chatter on the command.

\begin{table}[tbp]
\caption{G4 hover limits.}\label{tab:val_g4_lim}
\centering
\small
\begin{tabularx}{\textwidth}{@{}l l l l@{}}
\toprule
Metric & Quiet & Typical & Harsh \\
\midrule
R20 & 1~\% & 1.5~\% & 3~\% \\
C20 & 0.004 & 0.006 & 0.012 \\
R4 & 3~\% & 4~\% & 6~\% \\
Clip fraction & 0 & 0 & 0.1~\% \\
\bottomrule
\end{tabularx}
\end{table}

\Cref{tab:val_g4} gives the harsh-profile results, which are the worst case.

\begin{table}[tbp]
\caption{G4 results on the harsh profile.}\label{tab:val_g4}
\centering
\small
\begin{tabularx}{\textwidth}{@{}L{4.6cm} L{3.6cm} Y@{}}
\toprule
Quantity & Result & Limit \\
\midrule
Hover R20 / C20 / R4 / clip & 0.016~\% / 0.00025 / 0.041~\% / 0 & 3~\% / 0.012 / 6~\% / 0.1~\% \\
Cruise R20 / C20 / R4 & 0.023 to 0.024~\% / 0.00065 to 0.00067 / 0.271 to 0.295~\% & 4.5~\% / 0.018 / 6~\% \\
Command chatter, nearest to the flag & $4.12\times 10^{-4}$ & flag at $5\times 10^{-4}$ \\
Hover horizontal RMS / mean height error & 0.027 to 0.047~m / $-0.029$ to $+0.013$~m & 0.15~m / 0.1~m \\
PID guard at mode transitions & 54 transitions per key, none checked, none failed; largest INDI narrow-band RMS $3.46\times 10^{-4}$ & flag $5\times 10^{-4}$, failure above $1\times 10^{-3}$ \\
\bottomrule
\end{tabularx}
\end{table}

All 18 flights pass on both keys, with no key disagreement.
The guard compares INDI with PID at mode transitions where INDI shows a narrow line.
At the final configuration no transition reached the flag level, so none was compared.
The PID comparison exists only from an earlier firmware, where INDI and PID narrow-band values at the same transitions lay within a ratio of 0.78 to 1.14.

\indifigure{g4_noise_effort}{Rotor-speed ripple and motor-command RMS of the INDI controller for the quiet, typical and harsh noise profiles, in hover and cruise (three seeds each). Every run is far below its limit.}{fig:val_g4}

\begin{keyresult}[G4 verdict: \passed]
All 18 flights pass on both keys.
Evidence limits: harsh-cruise rate flags are informational; the PID comparison exists only for four harsh flights and comes from an earlier firmware; a 20~Hz take-off transient exists in both stacks at 0.05~\% of full command.
Item 3 (steady state) on cruise flights is not part of the G4 verdict: the scorer reports a miss on 9 of 9 cruise rows, in both stacks, because its window falls after the automatic landing starts.
\end{keyresult}

\section{G5: modes and transitions}\label{sec:val_g5}

G5 flies the INDI stack through all flight modes and transitions, three seeds per row, with a second key where a second script exists.
\Cref{tab:val_g5} gives the rows.
A row result is the worst of the three seeds.

\begin{table}[tbp]
\caption{G5 rows, limits and worst-case results.}\label{tab:val_g5}
\centering
\footnotesize
\begin{tabularx}{\textwidth}{@{}L{3.1cm} L{5.6cm} Y@{}}
\toprule
Row & Limit & Result \\
\midrule
Acro steps & RMS rate error at most 10~deg/s on 100~deg/s steps, no loss of control & 3 of 3, worst 1.94 \\
Stabilised release & tilt at most 3~deg, yaw rate at most 10~deg/s & 3 of 3, worst tilt 0.12 \\
Altitude, position steps & height overshoot at most 0.3~m; horizontal overshoot at most 0.5~m; settling at most 4~s & 3 of 3 each \\
Position hold in 5~m/s wind & 3-D RMS at most 0.3~m & 3 of 3, worst 0.062~m \\
Mission, 4 waypoints & within 1.0~m & 3 of 3, worst 0.970~m \\
Take-off, land & overshoot at most 0.5~m, edge step at most 10~\%, touchdown speed at most 0.8~m/s, tilt at most 10~deg & 6 of 6 \\
Return to launch & height error at most 1~m & 3 of 3, 0.045 to 0.062~m \\
Offboard trajectories o1 to o7 & horizontal RMS within $H_\mathrm{lim}$, vertical RMS at most 0.3~m & 21 of 21 under the amended limit (20 of 21 as first written) \\
Offboard velocity step & vx RMS at most 0.2, end vx at most 0.1, height at most 0.3~m & 3 of 3 \\
Mode switches & thrust step at most 20~\%, attitude jump at most 5~deg & 3 of 3, worst 3.98~deg \\
Failsafes in hover (RC, link, offboard loss) & hold entered within 1~s, excursion at most 1~m & 9 of 9, entry within 0.51~s, excursion 0.121~m \\
Failsafes at 3~m/s & same, excursion from the commanded setpoint & 9 of 9, excursion 0.16 to 0.25~m \\
Kill and re-arm & motors off within 1.25~ms of the kill; all INDI states zero; command change in the first second at most 5~\%; take-off step at most 10~\% & Met: 0.00 to 1.25~ms; zero; at most 0.01~\%; 4.0 to 4.7~\% \\
Stabilised full scale & no loss of control, tilt at most the mode limit plus 5~deg, attitude error at most 5~deg one second after the sticks centre & 3 of 3 \\
Estimator yaw reset & excess at most 1~deg over the no-event level, excursion at most 0.5~m & 3 of 3 \\
Estimator horizontal reset & true excursion at most 0.5~m, no oscillation & 3 of 3, 0.234 to 0.248~m \\
\bottomrule
\end{tabularx}
\end{table}

The offboard trajectory limit is $H_\mathrm{lim}=\max(0.5\,\mathrm{m},\ \tau_\mathrm{att} v_\mathrm{rms})$.
Here $\tau_\mathrm{att}=1/\param{MC_ROLL_P}=0.125$~s and $v_\mathrm{rms}$ is the RMS of the reference horizontal speed.
For trajectory o2 the limit is 0.5685~m and INDI reaches 0.5027~m.
Against a flat 0.5~m this one value would fail by 0.0027~m.
The margins to $H_\mathrm{lim}$ of the other trajectories range from 0.043 to 0.456~m.
On the failsafe test at 3~m/s, the PID reference has an excursion of 0.46 to 0.51~m on seed~1, against 0.16 to 0.25~m for INDI.

\indifigure{g5_modes}{Worst transient of each mode and transition test, as a percentage of its acceptance limit (INDI, three seeds). All rows are within their limits; the mission waypoint error (0.97~m of 1.0~m) and offboard trajectory o2 (0.5027~m of 0.5685~m) are the thinnest.}{fig:val_g5}

\begin{keyresult}[G5 verdict: \passed]
All rows that were tested pass.
Evidence limits:
\begin{itemize}
  \item Mission and return have one scoring script only; the mission margin to its limit is 0.03~m.
  \item Not tested, by owner decision: a false maybe-landed flag, a vertical estimator reset, acceleration-only offboard setpoints, and take-off on an offboard setpoint alone.
  \item The saturation part of the stabilised full-scale test was dropped, because no stack drives four motors above 0.99 for 1~s in this test (INDI reaches 0.869). Saturation under the position loop is covered by G3 corner C3, which INDI fails.
  \item Some second-key scripts print a different verdict on kill, re-arm and estimator-reset rows because they use older reference definitions. Read against the gate definitions, all values are within limits.
\end{itemize}
\end{keyresult}

\section{G6: defaults and identification}\label{sec:val_g6}

G6 asks whether the identification chain that ships with the controller can start from wrong physical parameters and still produce a controller that passes G2, G3 and G5.
The wrong starts are twice and half the true mass, inertia, motor time constant and thrust constant, one at a time (8 starts) plus the all-high and all-low corners.
The chain runs once, the identified values are loaded, and nothing else is edited.
G6 therefore requires G3.

One further test is flown alone: item 7.
It flies the 60 corner cases of G3 un-identified, with the identification flag off, controller gains halved and the actuator time constant at 0.05~s.
The pass rule is no loss of control in any valid run.
The owner reduced this test from about 100 Latin-hypercube runs plus corners to the 60 corners.

\begin{table}[tbp]
\caption{G6 results.}\label{tab:val_g6}
\centering
\small
\begin{tabularx}{\textwidth}{@{}L{4.8cm} Y@{}}
\toprule
Test & Result \\
\midrule
Requires G3 & G3 is final NOT PASSED, so G6 cannot pass \\
Item 7, no loss of control & Fails: 6 losses of control in 58 runs. Two runs lost control; four more were stopped by the flight watchdog (height loss, or tilt above 60~deg) and count as losses \\
Invalid runs (harness) & 2, not flown again because they cannot change the verdict \\
Reported only (key 1 / key 2 fail counts, 54 scored runs) & Oscillation over the whole run: 21 / 22. Steady state: 2 / 2. Tracking: 18 / 18. Recovery: 30 / 30. Rate mode: 0 / 0. Overall pass: 12 of 54 on both keys \\
Time-constant identification & 0 of 30 cells failed in a simulation grid; time constants within 3.7~\%, delay offset about $+1.75$~ms. One airframe, mostly one seed \\
Wrong-start identification (10 starts) & Not flown \\
Other airframes (about 1~kg, above 5~kg) & Not flown \\
\bottomrule
\end{tabularx}
\end{table}

The oscillation count covers the whole run.
Item 7 asks for it only inside the identification manoeuvre, and that window score was not made.
One run was scored before the owner changed the recovery rule.
The failures cannot be split by corner, because the run list does not record which corner each run flew.

\begin{keyresult}[G6 verdict: \notpassed]
G3 is not passed, and the un-identified start loses control in 6 of 58 runs.
The identification procedure and its defaults stay in the report as the method; \cref{sec:lim_g6} lists what was not covered.
\end{keyresult}

\section{G7: code}\label{sec:val_g7}

G7 tests the firmware as code.
\Cref{tab:val_g7} gives each check and its result at the final head of the firmware-only branch.

\begin{table}[tbp]
\caption{G7 checks.}\label{tab:val_g7}
\centering
\small
\begin{tabularx}{\textwidth}{@{}l L{4.2cm} Y@{}}
\toprule
Check & Test & Result \\
\midrule
(a) & Stock flight-control modules are unchanged (zero diff against the main line) & Pass: 0 lines over the six stock paths, over 75 files \\
(b) & Unit tests pass, and each behaviour test fails when the behaviour is broken & Pass: 12 of 12 test programs, 239 tests (244 after the arming check was added). A mutation audit of 66 deliberate code breaks killed all 9 survivors. Seven weak or missing test rows were closed by new tests \\
(c) & Firmware default parameters equal the validated parameter files & Met: 0 mismatches in gains. Physical values (maximum rotor speed, actuator time constant) are per airframe \\
(d) & Code format check is clean & Pass: 53 files clean \\
(e) & Flash and RAM measured on one NuttX board from the build & Pass on the v5x board with both modules: $+64\,792$~B flash at first (8\,964~B left), $+66.1$~kB after the arming check (7.7~kB left). The v6x board overflows by 23.5~kB \\
(f) & Instruction files and naming rules followed; whole-branch review & Pass: delta review with 0 blockers and 0 major findings; 92 review items closed, justified or superseded \\
(g) & Firmware-only branch from the main line builds and has no project-only path & Pass: simulation target builds, zero project identifiers \\
\bottomrule
\end{tabularx}
\end{table}

Some code is justified rather than unit-tested.
The wiring of some parameters into the library is covered by the simulation rows of G2 and G5.
A permanently saturated motor with frozen telemetry is not detected, by design, and a test pins that behaviour.
Part of the position module (a position-and-velocity gain cascade) is not reachable, because the module always runs acceleration-only; it is listed as an open item.

\begin{keyresult}[G7 verdict: \passed, with justified limits]
Every check passes.
The firmware does not fit the v6x board with both modules.
Production boards do not build the modules; only the simulation target does.
\Cref{sec:lim_impl} lists the open items.
\end{keyresult}

\section{G8: report}\label{sec:val_g8}

G8 is this report: seven mechanical checks of its content pass, and the unverified claims it carries are marked in the text and counted.

\section{Amendments to the tests after the first run}\label{sec:val_amendments}

The owner amended these tests after their first run.
Each amendment fixes a defect of the test.

\begin{description}[leftmargin=1.4em,labelindent=0pt]
  \item[G1.1] Logger completeness replaced the original gap rule, which simulator holes broke; the mode-change spread and latency limits were set to 10 and 20~ms.
  \item[G1.5, G5 events] A discrete message event (stream stop, kill, estimator reset) is timed from the firmware's own applied flag, so that the measured delay does not include the message transport.
  \item[G5.6, G5.11] The false maybe-landed test and the vertical estimator reset are not tested, because the simulation cannot inject them: the landing detector needs ground contact, and the estimator accepts only a horizontal external reset.
  \item[G5.8] Failsafe tests at 3~m/s measure the excursion from the setpoint the failsafe commands, and the link-loss timeout is timed from the firmware's own flag.
  \item[G5.9] ``Motors off within one command period'' is timed from the kill flag, not from the commander request.
  \item[G5.2] The height settling band is the larger of 2~\% and 0.10~m, because the estimator height error is that large.
  \item[G5.10] The full-scale test is flown in stabilised mode with the tilt limit of that mode plus 5~deg, and the saturation part is dropped unless four motors exceed 0.99 for one second in both stacks (they do not).
  \item[G5.11] The horizontal reset is judged on the true excursion of at most 0.5~m and the tilt jump is only reported.
  \item[G5.5] The horizontal limit for moving trajectories is $\max(0.5\,\mathrm{m},\ \tau_\mathrm{att}v_\mathrm{rms})$, because the flat 0.5~m limit is a steady hover limit and does not fit a moving trajectory.
  \item[G2] Where $|L|>1$ within 10~deg of $-180^\circ$, the gain margin is the gain reduction $20\log_{10}|L|$, which must be at least 6~dB, because the plateau reading misses such crossings.
  \item[G4 chatter] A 20 to 100~Hz line counts as chatter only when it is 20 times above the median and as large as one command step ($5\times 10^{-4}$).
  \item[G4 and G3 windows] Detector windows that overlap the 3~s after a mode change, or the landing, are skipped, because mode transitions are judged by G5.
  \item[G4 guard] Skipped windows are checked against PID: the row fails if INDI shows a line that PID does not, or an INDI line above $10^{-3}$.
  \item[G3 factors] The controller-versus-plant gain error was narrowed from 0.5 to 2 times to 0.8 to 1.2 times, because the identification step is expected to be within 10~\%; the motor up and down time constants both lie in 0.5 to 3 times nominal; the estimator GNSS delay mismatch was reduced from $\pm 100$~ms to $\pm 50$~ms, because a correct bring-up leaves less.
  \item[G3 controller time constant] The controller's actuator time constant is drawn like an identification error around the plant value, instead of being fixed.
  \item[G3 items 3 to 6] There is no pass quota for them: every miss in a valid run is a failure.
  \item[G3 harness aborts] A harness abort is an invalid run, flown again, and never a pass.
  \item[G3 steady state, events] The steady-state window is the last 10~s of a hold of at least 20~s, and faults and gusts happen in a separate event hover before the main segment, so that the recovery and steady-state windows do not overlap.
  \item[G3 oscillation windows] Item 2 skips every 4~s window that contains a fault or gust start.
  \item[G3 item 5] Recovery is judged against a static reference, and, after the owner's reading, requires both a recovery within 5~s and a clean 5 to 10~s afterwards.
  \item[G3 item 4] A row where PID also fails the same draw is a test defect and is not counted, if INDI is not worse.
  \item[G6 item 7] It was added after the first robustness run, flown with the identification flag off, and reduced to the 60 corners.
\end{description}

Most amendments fix timing, windows or references.
Two change a limit or a range: the G5.5 limit now grows with trajectory speed, and the G3 gain-error box was narrowed.
Both are written above with their reason.
